\section{The Christian Social Order \& Pythagoras}

Christopher Ferrara gives a very basic, but nonetheless interesting tour of Western civilization here\footnote{\url{http://www.youtube.com/watch?v=3P9nFumxKp8}}. I hadn't heard of Werner Jaeger\footnote{\url{http://en.wikipedia.org/wiki/Werner\_Jaeger}}, but I've ordered his three volume set. One sentence really springs to mind: \emph{``The State was the atmosphere within which the soul breathed''}.

This was the common link or identical conception that joined Hellenism with Christian Europe. What it seems to be saying is that the State and the soul \emph{are equally spiritual, that they in fact co-arise}, and that you cannot oppose them against one another. Both have a transcendent origin and divine destiny. However, experience shows us that they do (in fact) conflict. Witness the great Greek tragedy of Antigone\footnote{\url{http://en.wikipedia.org/wiki/Antigone}}, in which the Greek mind pitched this tension to a great height. So intellectually, they have to be co-equals or resolved in a higher One-in-Many, but experientially they run into conflicts, as Athens discovered during the brutal struggle with Sparta, when she ``destroyed the village in order to save it''\footnote{\url{http://en.wikipedia.org/wiki/Melian\_dialogue}}.

Ferrara argues that the role or place of the Church of St. Peter (we would say the external, exoteric Church) was actually to arbitrate this tension, to decide in questions of natural or moral law when they come into conflict, by using the supernatural mind of Christ. Regardless of what you think this caused in history, this is a brilliant and sound way of explaining the evolution of the Hellenic and Celtic world (which shared these tensions) into medieval and traditional Europe. The Church provides the balance, and mediates between ``Antigone'' and the divine state. What this proves is that the Soul and the State are relatively divine: that is Spirit exists (as well as celestial powers) that transcend either. However, when the Soul (or the State as a ``soul'' in the divine body of the King\footnote{\url{http://www.libraryofsocialscience.com/ideologies/docs/the-kings-two-bodies/index.html}}) is raised by grace to a high pitch of holiness, you can see exoteric conflict at the terrestial level of common law or State rule, or perhaps between two noble souls. So the conflict is both necessary and superficial (if both sides are guided by the celestials, angels and God), \& it can be resolved by God, provided both parties adhere to the Spirit.

There is no question that this process broke down in Europe, and we began to see irresolvable conflicts between holy souls, with Christian Germany drawing the sword\footnote{\url{http://wwi.lib.byu.edu/index.php/Wilhelm\_II\%27s\_War\_Speeches}} against Christian France, both sides singing Christmas hymns in the trenches. The secular narrative states that this only proves that we had reached a dead end, and that in fact the entire narrative even back to Socrates or ancient Israel was in fact a gigantic fraud, and that there is no such thing as positive revealed religion, transcendent souls, or a Church (although various parties within modern Western civilization keep various elements of these -- the classical liberals keep the transcendent soul, the evangelicals still believe in a Church, etc.). But the full synthesis was shattered, beginning (officially) with John Locke and our own Founding Fathers, who accepted the shattering of the synthesis begun by Luther and others within even the Church (read Luther on the scholastics).

So we have a second tension that develops : what happens when the Church is no longer able to arbitrate between Antigone and the folk-State or the semi-divine \emph{polis}? Again, the modern answer is philosophical despair, and the flight into comfort, democracy, and a positive law that depends only upon the decree of Hobbes' Leviathan, or a Machiavellian prince. In other words, we have Chaos, as the older elements never completely retreat, but actually help prop up the mass of confusion and allow it to continue. The whole rickety affair, like the Old One Hoss Shay\footnote{\url{http://www.legallanguage.com/resources/poems/onehossshay/}}, holds together against all odds until it just doesn't any more.

The esotericist may well be able to provide some help here, as the Church of John has come into more light thanks to the plight of the external Church of Peter. The alchemist or the esotericist may help to heal the conflicts between various elements within the West, but it is important to note that healing will (by Logos) have to occur along the lines of the existing synthesis, shattered as it was and is. Thus there is an implicit endorsement of General Law by Mouravieff, Pythagoras, and other perennial thinkers (as we have noted on this blog). That is, the esotericist attitude is always that it is best for the perverse man to die, as there is no help for such a one in this plane of existence (sayings of Pythagoras). Thus, Romans 13 and the ``sword'' of the State, along with the Ten Commandments, the Tao, and the necessity of folk and moral traditions still stands. To give an example, after the Pythagoreans were assassinated in a Greek city-state (only 2 escaping the burning house) they withdrew from politics and left the wicked to themselves. This is a sound attitude, and it is also a sound attitude for Mouravieff to point out that Chaos can sometimes reach a point in which it is preferable for the good man to favor Law over this Chaos, a principle also taught by Pythagoras: ``strengthen those of the Law, do not help those who are against it''. This is the fundamental anti-Revolutionary attitude of even the Church of John. ``Let the dead bury the dead''. \emph{Christ did not come in order to give revolutionaries a chance to create utopia on earth by agitating for liberty from the Logos}. See Groen van Prinsterer\footnote{\url{http://en.wikipedia.org/wiki/Guillaume\_Groen\_van\_Prinsterer}}`s work, for even a Reformed rejection of Revolution. Or, as James puts it, Love the Brotherhood, Fear God, Honor the King. There is Order in three sentences.

Thus, we see that classical liberalism may be the highest secular personal ethos to ever reflect Divine Truth, but it is (correspondingly) \emph{the absolutely most insidious form of substitute or ersatz for either civic religion or political General Law}. Now why would someone say this? \textbf{This is because it makes a Law out of that precisely which is supposed to transcend Law, to exist outside of the Law, and to fulfill or supersede the Law.} Isn't this sounding almost Phariseeical? This supercession is in reality and truth, not in political or material matters, because the two must always be separate (Christ and Mammon).

Now from the Left, this is precisely what one should not do; instead one must mingle Mammon and Christ, because one really doesn't believe in either, but rather in the sorcerer-state illusion that is created by the co-mingling. Here is Culianu defending precisely\footnote{\url{http://books.google.com/books?id=DDVkNj97dMoC&pg=PA105&lpg=PA105&dq=sorcerer+state&source=bl&ots=RUiO2pcH\_q&sig=EQZguDlKDvyVy68VREa8SLUY6Hg&hl=en&sa=X&ei=zxzxUr\_pMaTKsQTqoYF4&ved=0CEUQ6AEwAw\#v=onepage&q=sorcerer\%20state&f=false}} this Leftist position. Since an irresolvable hypocrisy is present and presents itself in the Leftist mind, it leads to a pre-emptive charge of hypocrisy against truly traditional orders, which invariably present (to the modern mind) the spectacle of ``good men'' winking at ``evil'' by endorsing such things as slavery, genocide, repression, prejudice, etc. Of course, these are loaded terms, since from a higher point of view, the Right is not endorsing them at all, they are destroying them by ignoring them, which is the only way to do away with them (contra the Utopians). Thus, \emph{render unto Caesar, what is Caesar's}, etc. And, one may add, let Caesar render unto God what is God's. This looks like hypocrisy to the fundamentalistic mind of the modern secular classical liberal, because they have made a ``Law'' out of ``Grace''. Like the Pharisees, they are angry that Christ comes eating and drinking, endorsing Caesar, and mingling with sinners (in today's narrative, those who are not ``\emph{politically correct}''). To them, it sure looks like Christ should have put the Romans in their place; they are often fond of comparing conservatives to the hypocrite Pharisees, but these days, classical liberalism looks more like the civic religion of choice for America, replete with its own sacraments, high priest, martyrs, and accoutrements that justify a religious lineage. The greatest political or moral sins of our day are (in this order) bigotry, intolerance, and prejudice, all of which are viewed as fundamentally hypocritical.

The \emph{veredictum} of Tradition lies against this conclusion. Pythagoras had no problem endorsing law-giving tyrants, nor did Plato, Mouravieff, and apparently Christ himself, who came to bring a sword to the earth and to make the Ten Commandments more rigorous. None of these endorsed tyrants were Leftists, but were emphatically prototypes of ``the Right'' (at least from a Leftist perspective). They were corrupt men, but men of Order, especially by comparison with today. If I were to look hard enough, I'll bet I could find some reference in Tomberg to such matters, as well. Someone told me that there can't be too much a separation between politics and spirituality, or else ``it doesn't make much sense anymore''. Well, men of Tradition can always simply withdraw from politics (like Ernst Junger or Pythagoras) and turn a blind eye to the inevitable backlash against Revolution. This (of course) will also be branded as hypocrisy.

Actually, there should be a huge chasm between politics and spirituality, because ``most men'' shouldn't dabble in politics, and they should follow a religion, rather than worry about being ``spiritual''. So perhaps neither should exist: the biggest gap possible! The State and the individual, however, will continue to share a transcendent purpose that neither can deny to the other, as long as higher religious vehicles and esoteric practice fecundate society and ritual with what TS Eliot calls the necessity\footnote{\url{http://www.theimaginativeconservative.org/2012/06/t-s-eliots-christianity-and-culture.html}} of a ``Christian culture''. The problem isn't disjunction between religion and state, it's that the Left only sees a particular kind of conjunction which they object to. Their own version (a soft tyranny formed by democratic and corporate control of thought\footnote{\url{http://www.amazon.com/Democracy-Incorporated-Managed-Inverted-Totalitarianism/dp/069114589X}}) seems not merely benign but positively self-evident. Whatever Order manifests itself as predominant, it will have the same enemies that Christianity has today -- or, at least, they will be saying the same things.

It's important, in the moral confusion and intellectual Chaos, to remember what men of Tradition (including the very classical liberals who progenitated our latter day mutation) said concerning the nature of the State, the nature of the soul, the nature of God. To do so is the only kind of revolutionary act that has any truth in it at all.

I think there is more to say on the subject, and much farther to see, by those with clearer eyes and hearts and heads than my own. Whatever they find, I will wager that there is a kind of frightening confusion of all categories that occurs within the Leftist dialectic, and allows it to project its own hypocrisy as a form of shadow upon those who see that some evils in the world are not capable of eradication, least of all by overthrowing what little Order man is capable of half way keeping up.


\flrightit{Posted on 2014-02-04 by Logres}

\begin{center}* * *\end{center}

\begin{footnotesize}
\begin{sffamily}
\texttt{Logres on 2014-02-04 at 22:55 said:}

Should have included: ``Liberalism. . . tends to release energy rather than accumulate it, to relax, rather than to fortify. It is a movement not so much defined by its end, as by its starting point; away from, rather than towards, something definite. Our point of departure is more real to us than our destination; and the destination is likely to present a very different picture when arrived at, from the vaguer image formed in imagination. By destroying traditional social habits of the people, by dissolving their natural collective consciousness into individual constituents, by licensing the opinions of the most foolish, by substituting instruction for education, by encouraging cleverness rather than wisdom, the upstart rather than the qualified, by fostering a notion of getting on to which the alternative is a hopeless apathy, Liberalism can prepare the way for that which is its own negative: the artificial, mechanized or brutalised control which is a desperate remedy for its chaos. (CC, 12)'' TS Eliot

\hfill

\texttt{Caleb Cooper on 2014-02-06 at 14:41 said:}

Tomberg on Revolution: ``... it is not the difference in dogma nor that of cult or ritual which is the cause of Cain's crime , but uniquely the \emph{pretension to equality} italics in original or, if one prefers, the \emph{negation of hierarchy.} Here also is the world's first revolution -- the archetype of all revolutions... For the cause of all wars and revolutions -- in a word, of all violence -- is always the same: the negation of hierarchy.

On authority, law, and their dependence on the office of Emperor: ``This is why the first idea that the Emperor Card naturally evokes is that of the \emph{authority} underlying \emph{law}. The thesis which proceeds from the meditation on the three preceding Arcana is that all authority has its source in the ineffable divine name of YHVH and that all law derives from this.

This is why authors of the Middle Ages could not imagine Christianity without an Emperor, just as they could not imagine the Universal Church without a pope. Because if the world is governed hierarchically, Christianity or the \emph{Sanctum Imperium} cannot be otherwise. Hierarchy is a pyramid which exists only when it is complete. And it is the Emperor who is at its summit. Then comes the kings, dukes, noblemen, citizens and peasants. But it is the crown of the Emperor which confers royalty to the royal crowns from which the ducal crowns and all the other crowns in turn derive their authority.

Europe is haunted by the shadow of the Emperor... because the emptiness of the wound \emph{speaks}... ''

\hfill

\texttt{Logres on 2014-02-08 at 00:41 said:}

Found in Tomberg: ``The crime of all revolutionaries is emotional impatience''... 

\hfill

\end{sffamily}
\end{footnotesize}

